% Part: propositional-logic
% Chapter: propositional-logic
% Section: preliminaries

\documentclass[../../../include/open-logic-section]{subfiles}

\begin{document}

\olfileid{pl}{syn}{pre}

\olsection{প্রাথমিক ফলাফল}

\begin{thm}[\emph{!!{formula}s-এর উপর আরোহের নীতি}]
  \ollabel{thm:induction}
  কোনো ধর্ম~$P$ যদি সব পরমাণু !!{formula}s-এর জন্য সত্য হয় এবং
  এমন হয় যে
  \begin{enumerate}
    \tagitem{prvNot}{$!A$-এর জন্য ধর্মটি সত্য হলেই $\lnot !A$-এর
      জন্যও সত্য হয়;}{}
    \tagitem{prvAnd}{$!A$ ও~$!B$-এর জন্য সত্য হলেই
      $(!A \land !B)$-এর জন্যও সত্য হয়;}{}
    \tagitem{prvOr}{$!A$ ও~$!B$-এর জন্য সত্য হলেই
      $(!A \lor !B)$-এর জন্যও সত্য হয়;}{}
    \tagitem{prvIf}{$!A$ ও~$!B$-এর জন্য সত্য হলেই
      $(!A \lif !B)$-এর জন্যও সত্য হয়;}{}
    \tagitem{prvIff}{$!A$ ও~$!B$-এর জন্য সত্য হলেই
      $(!A \liff !B)$-এর জন্যও সত্য হয়;}{}
  \end{enumerate}
  তবে সব !!{formula}s-এর জন্য $P$ সত্য।
\end{thm}

\begin{proof}
  ধর্ম~$P$-বিশিষ্ট সব !!{formula}s-এর সংগ্রহকে $S$ ধরি।
  স্পষ্টত $S \subseteq \Frm[L_0]$। $S$,
  \olref[fml]{defn:formulas}-এর সব শর্ত পূরণ করে: এতে সব পরমাণু
  !!{formula}s আছে এবং এটি !!{operator}s-এর অধীনে বদ্ধ।
  $\Frm[L_0]$ এই রকম ক্ষুদ্রতম শ্রেণি, তাই $\Frm[L_0] \subseteq S$।
  অতএব $\Frm[L_0] = S$, এবং প্রতিটি সূত্রের ধর্ম~$P$ আছে।
\end{proof}

\begin{prop}\ollabel{prop:balanced}
   $\Frm[L_0]$-এর যেকোনো !!{formula} \emph{সুষম}; অর্থাৎ তাতে
   বাঁদিকের ও ডানদিকের বন্ধনীর সংখ্যা সমান।
\end{prop}

\begin{prob}
\olref[pl][syn][pre]{prop:balanced} প্রমাণ করো।
\end{prob}

\begin{prop} \ollabel{prop:noinit}
কোনো !!a{formula}-এর প্রকৃত প্রারম্ভিক খণ্ড !!a{formula} নয়।
\end{prop}

\begin{prob}
  \olref[pl][syn][pre]{prop:noinit} প্রমাণ করো।
\end{prob}

\begin{prop}[একক পাঠযোগ্যতা]
$ \Frm[L_0]$-এর যেকোনো !!{formula}~$!A$-এর নিচের রূপগুলির ঠিক একটিতে
একটিমাত্র গঠনবিশ্লেষণ আছে:
\begin{enumerate}
\tagitem{prvFalse}{$\lfalse$।}{}

\tagitem{prvTrue}{$\ltrue$।}{}

\item কোনো $\Obj p_n \in  \PVar$-এর জন্য $\Obj p_n$।

\tagitem{prvNot}{কোনো !!{formula}~$!B$-এর জন্য $\lnot !B$।}{}

\tagitem{prvAnd}{কোনো !!{formula}s $!B$ ও~$!C$-এর জন্য
$(!B \land !C)$।}{}

\tagitem{prvOr}{কোনো !!{formula}s $!B$ ও~$!C$-এর জন্য
$(!B \lor !C)$।}{}

\tagitem{prvIf}{কোনো !!{formula}s $!B$ ও~$!C$-এর জন্য
$(!B \lif !C)$।}{}

\tagitem{prvIff}{কোনো !!{formula}s $!B$ ও~$!C$-এর জন্য
$(!B \liff !C)$।}{}
\end{enumerate}
তদুপরি, এই গঠনবিশ্লেষণটি \emph{একক}।
\end{prop}

\begin{proof}
$!A$-এর উপর আরোহ প্রয়োগ করি। যেমন, ধরি $!A$-কে $(!B \lif !C)$ ও
$(!B' \lif !C')$ হিসেবে পাঠ করার দুটি স্বতন্ত্র উপায় আছে। তাহলে $!B$ ও
$!B'$ অবশ্যই অভিন্ন (নইলে একটি অন্যটির প্রকৃত প্রারম্ভিক খণ্ড হতো)। তাই
$!A$-এর পাঠদুটি স্বতন্ত্র হলে তার কারণ অবশ্যই এই যে $!C$ ও $!C'$ একই
প্রতীকক্রমের স্বতন্ত্র পাঠ; আরোহ-অনুমান অনুসারে যা অসম্ভব।
\end{proof}


\begin{defn}[সমরূপ প্রতিস্থাপন]
যদি $!A$ ও $!B$ !!{formula}s হয় এবং $\Obj p_i$ একটি !!{propositional
variable} হয়, তবে $!A$-তে $\Obj p_i$-এর প্রতিটি উপস্থিতির জায়গায়
$!B$-এর একটি উপস্থিতি বসালে যে ফল পাওয়া যায়, তাকে
$\Subst{!A}{!B}{\Obj p_i}$ নির্দেশ করে। একইভাবে, $\Obj p_1$, \dots,~$\Obj p_n$-এর
স্থানে একই সঙ্গে !!{formula}s $!B_1$, \dots,~$!B_n$ বসানোর প্রতিস্থাপনকে
$\SSubst{!A}{\subst{!B_1}{\Obj p_1},\dots,\subst{!B_n}{\Obj p_n}}$
দিয়ে নির্দেশ করা হয়।
\end{defn}

\begin{prob} নিচের পাঁচটি !!{formula}s-এর প্রত্যেকটির জন্য নির্ধারণ করো,
 !!{formula}-টিকে \( \Subst{!A}{!B}{\Obj p_i} \) প্রতিস্থাপনের রূপে
 লেখা যায় কি না, যেখানে \( !A \) হলো (i) \( \Obj p_0 \); (ii)
 \( ( \lnot \Obj p_0 \land \Obj p_1) \); এবং (iii)
 \( ( ( \lnot \Obj p_0 \lif \Obj p_1 ) \land \Obj p_2 ) \)।
 প্রতিটি ক্ষেত্রে প্রাসঙ্গিক প্রতিস্থাপনটি নির্দিষ্ট করো।
  \begin{enumerate}
    \item \( \Obj p_1 \)
    \item \( ( \lnot \Obj p_0 \land \Obj p_0 ) \)
    \item \( ( ( \Obj p_0 \lor \Obj p_1 ) \land \Obj p_2 )  \)
    \item \( \lnot ( ( \Obj p_0 \lif \Obj p_1 ) \land \Obj p_2 ) \)
    \item \( (( \lnot ( \Obj p_0 \lif \Obj p_1 ) \lif ( \Obj p_0 \lor \Obj p_1 )) \land \lnot ( \Obj p_0 \land \Obj p_1 )) \)
  \end{enumerate}
\end{prob}

\begin{prob}
আরোহের সাহায্যে $\Subst{!A}{!B}{p}$-এর গাণিতিকভাবে কঠোর একটি সংজ্ঞা দাও।
\end{prob}

\end{document}
